\par\medskip
\begin{center}\hypertarget{native-164-article118}{}\pdfbookmark[2]{118.}{article-118}118.\end{center}
Si corporis coelestis distantia a terra alicui observationi respondens proxime iam nota est, hanc ab effectu \emph{aberrationis} liberare licet pluribus modis, qui methodis diversis in art. 71 traditis innituntur. Sit $t$ tempus verum observationis; $\theta$ intervallum temporis, intra quod lumen a corpore coelesti ad terram descendit, quod prodit ducendo $493^{\mathrm{s}}$ in distantiam; $l$ locus observatus, $l'$ idem locus adiumento motus geocentrici diurni ad tempus $t+\theta$ reductus; \SourceLeaf{165}$l''$ locus $l$ ab ea aberrationis parte purgatus, quae planetis cum fixis communis est; $L$ locus terrae verus tempori $t$ respondens (i. e. tabularis $20''{,}25$ auctus): denique ${}'L$ locus terrae verus tempori $t-\theta$ respondens. His ita factis erit
\par I. $l$ locus verus corporis coelestis ex ${}'L$ visus tempore $t-\theta$
\par II. $l'$ locus verus corporis coelestis ex $L$ visus tempore $t$
\par III. $l''$ locus verus corporis coelestis ex $L$ visus tempore $t-\theta$.
\par\noindent Per methodum I itaque locus observatus immutatus retinetur, pro tempore vero autem fictum $t-\theta$ substituitur, loco terrae pro eodem computato; methodus II soli observationi mutationem applicat, quae autem praeter distantiam insuper motum diurnum requirit; in methodo III observatio correctionem patitur a distantia non pendentem, pro tempore vero fictum $t-\theta$ substituitur, sed retento loco terrae tempori vero respondente. Ex his methodis prima longe commodissima est, quoties distantia eatenus iam nota est, ut reductio temporis $\theta$ praecisione sufficiente computari possit.

Quodsi autem haec distantia penitus adhuc incognita est, nulla harum methodorum immediate applicari potest: in prima scilicet habetur quidem corporis coelestis locus geocentricus, sed desideratur tempus et positio terrae a distantia incognita pendentia; in secunda e contrario adsunt haec, deest ille; denique in tertia habetur locus geocentricus corporis coelestis atque positio terrae, sed tempus deest cum illis datis iungendum.

Quid faciendum est itaque in problemate nostro, si in tali casu solutio respectu aberrationis quoque exacta postulatur? Simplicissimum utique est, orbitam primo neglecta aberratione determinare, quae quum effectum considerabilem nunquam producere possit, distantiae hinc ea certe praecisione demanabunt, ut iam observationes per aliquam methodorum modo expositarum ab aberratione purgare, orbitaeque determinationem accuratius iterare liceat. Iam in hocce negotio methodus tertia ceteris longe praeferenda erit: in methodo enim prima omnes operationes a positione terrae pendentes ab ovo rursus inchoandae sunt: in secunda (quae ne applicabilis quidem est, nisi tanta observationum copia adsit, ut motus diurnus inde elici possit) omnes operationes a loco geocentrico corporis coelestis pendentes denuo instituere oportet: contra in tertia (siquidem iam calculus primus superstructus fuerat locis geocentricis ab aberratione fixarum purgatis) omnes operationes praeliminares a \SourceLeaf{166}positione terrae et loco geocentrico corporis coelestis pendentes in computo novo invariatae retineri poterunt. Quin adeo hoc modo primo statim calculo aberrationem complecti licebit, si methodus ad determinationem orbitae adhibita ita comparata est, ut valores distantiarum prodeant prius, quam tempora correcta in calculum introducere opus fuerit. Tunc aberrationis quidem caussa calculus duplex haud necessarius erit, uti in tractatione ampliori problematis nostri clarius apparebit.
